%% ---------------------------------------------------------------------------
\section{Signing and verifying where resources are short}
\label{sec:edge}

A camera, a body-worn recorder or a field encoder signs on something closer to
a Raspberry Pi~5 than to a laptop, and a doorbell or a sensor on something
smaller still. This section measures both. The Pi is an application processor,
not a secure element, and Section~\ref{sec:constrained} says what that leaves
open.

\subsection{Signing an encoded stream}
\label{sec:edge-video}

Cameras and encoders bind the encoded bitstream, not raw pixels, and C2PA
binds video through hashes of ISO BMFF segments rather than per-frame
signatures. C2PA~2.5 gives two methods for a live stream~\cite{c2pa25}. In the
first, every segment carries a manifest of its own, with a BMFF hash of the
segment, a live-video assertion carrying a sequence number and stream
identifier, and a signature with the signer's chain. In the second, the
stream's first manifest carries the chain and a session key bound to the
signer, and each segment carries only a small Verifiable Segment Info map,
its sequence number, hash and manifest identifier, signed by the session key.
We signed real streams both ways on one core of the Pi~5: \NstreamSeconds{} seconds of 1080p30 H.264 at 12\,Mbit/s and of 4K30 HEVC at 40\,Mbit/s, in fragmented MP4 with one fragment per second, so that a segment of $k$ seconds is $k$ fragments and even six-second segments give \NstreamSixSegs{} per pass. As an upper bound we also signed every encoded frame with the
chain attached, which C2PA does not define.

\begin{table*}[t]
\caption{Signing an encoded stream on the Raspberry Pi 5, one core. \emph{Frame} signs every encoded access unit with the chain attached (an upper bound, not a C2PA method); \emph{manifest} is C2PA~2.5 Section~15.4, a manifest per segment; \emph{VSI} is Section~15.5, the chain once and a session-key signature per segment. Provenance rate is the bitrate the signatures add; \emph{latency} is from segment close to signature (last fragment's hash plus signing). A player can authenticate a segment only once it holds all of it, so verification lags by at least the segment length. The provenance rate depends on the size of each signature and on how often one is sent, not on the video, so it is the same for both streams; the latency is not. Streams of 288\,s; medians over every segment of three passes.}
\label{tab:stream}
\small
\setlength{\tabcolsep}{3.2pt}
\begin{tabular}{llrrrrrrrrrrr}
\toprule
 & & Seg. & \multicolumn{2}{c}{Ed25519} & \multicolumn{2}{c}{PS256} & \multicolumn{2}{c}{ML-DSA-44} & \multicolumn{2}{c}{ML-DSA-65} & \multicolumn{2}{c}{ML-DSA-87}\\
\cmidrule(lr){4-5}\cmidrule(lr){6-7}\cmidrule(lr){8-9}\cmidrule(lr){10-11}\cmidrule(lr){12-13}
Stream & Method & (s) & kbit/s & ms & kbit/s & ms & kbit/s & ms & kbit/s & ms & kbit/s & ms\\
\midrule
1080p H.264 & frame & -- & 413.3 & 0.10 & 676.8 & 3.59 & 2,739.8 & 0.26 & 3,687.1 & 0.37 & 4,943.3 & 0.48\\
 & manifest & 1 & 22.3 & 1.07 & 31.1 & 4.56 & 99.9 & 1.22 & 131.5 & 1.32 & 173.3 & 1.44\\
 & manifest & 2 & 11.2 & 1.08 & 15.6 & 4.58 & 50.0 & 1.24 & 65.7 & 1.35 & 86.7 & 1.44\\
 & manifest & 4 & 5.6 & 1.06 & 7.8 & 4.56 & 25.0 & 1.22 & 32.9 & 1.33 & 43.3 & 1.44\\
 & VSI & 1 & 2.3 & 1.05 & 3.8 & 4.56 & 21.1 & 1.18 & 28.3 & 1.28 & 38.8 & 1.37\\
 & VSI & 2 & 1.1 & 1.06 & 1.9 & 4.57 & 10.6 & 1.20 & 14.1 & 1.29 & 19.4 & 1.39\\
 & VSI & 4 & 0.6 & 1.04 & 1.0 & 4.55 & 5.3 & 1.18 & 7.1 & 1.30 & 9.7 & 1.39\\
 & VSI & 6 & 0.4 & 1.09 & 0.6 & 4.60 & 3.5 & 1.21 & 4.7 & 1.30 & 6.5 & 1.43\\
\midrule
4K HEVC & frame & -- & 413.3 & 0.15 & 676.8 & 3.65 & 2,739.8 & 0.35 & 3,687.1 & 0.45 & 4,943.3 & 0.57\\
 & manifest & 1 & 22.3 & 3.35 & 31.1 & 6.86 & 99.9 & 3.53 & 131.5 & 3.64 & 173.3 & 3.75\\
 & manifest & 2 & 11.2 & 3.35 & 15.6 & 6.85 & 50.0 & 3.53 & 65.7 & 3.62 & 86.7 & 3.73\\
 & manifest & 4 & 5.6 & 3.33 & 7.8 & 6.84 & 25.0 & 3.49 & 32.9 & 3.62 & 43.3 & 3.74\\
 & VSI & 1 & 2.3 & 3.34 & 3.8 & 6.85 & 21.1 & 3.49 & 28.3 & 3.59 & 38.8 & 3.68\\
 & VSI & 2 & 1.1 & 3.33 & 1.9 & 6.84 & 10.6 & 3.47 & 14.1 & 3.57 & 19.4 & 3.67\\
 & VSI & 4 & 0.6 & 3.32 & 1.0 & 6.83 & 5.3 & 3.44 & 7.1 & 3.55 & 9.7 & 3.63\\
 & VSI & 6 & 0.4 & 3.34 & 0.6 & 6.86 & 3.5 & 3.46 & 4.7 & 3.58 & 6.5 & 3.65\\
\bottomrule
\end{tabular}
\end{table*}


Table~\ref{tab:stream} gives the result. Time is not the constraint. From the moment a segment closes, hashing its last fragment and signing take at most \NstreamLatMax\,ms, and the 99th percentile of the whole per-segment work, every fragment's hash and the provenance built and signed, is at most \NstreamPnnMax\,ms, against segments of one to six seconds. Counting all of that work, either C2PA method uses at most \NstreamCpuMax\% of one core; signing every frame uses up to \NstreamFrameCpuMax\%. These are feasibility figures from \NstreamSamplesMin{} or more segments per configuration, not a characterisation of tail latency under load. What the method decides is bandwidth. Under
ML-DSA-65 at 1080p, a manifest per one-second segment adds
\NstreamManOneKbps\,kbit/s, \NstreamManOneShare\% of the video; the session-key
method adds \NstreamVsiOneKbps\,kbit/s, \NstreamVsiOneShare\%, after a first
manifest of \NstreamVsiInit{} bytes; and signing every frame adds
\NstreamFrameOneKbps\,kbit/s, \NstreamFrameOneShare\% of the stream. Longer
segments divide the per-segment cost but delay authentication: a player can
check a segment only once it holds all of it, so verification lags capture by
at least the segment length. The session key here uses the signer's
algorithm, so the stream stays post-quantum end to end; a classical session
key under a post-quantum signer would shrink each segment's signature and
leave every segment forgeable by a quantum adversary.

A verifier that knows which stream it is playing replayed each stream with ten manipulations (Table~\ref{tab:streamdetect} in the supplement). Under both C2PA methods it detected a dropped, reordered, replayed or substituted segment, bytes truncated within a segment, and a segment taken from another stream by the same signer: truncation and substitution by the hash, the cross-stream segment by its stream identifier, and the rest by the sequence number. Without the initial manifest a session-key stream cannot be verified at all, and segments signed under an expired session key are rejected. One manipulation is not detected: removing the last segments. A stream that stops early looks like one whose tail was cut, unless its length is announced in signed data or the player gives up after a timeout; with an announced length, the cut is caught. A per-segment signature without a sequence number would also have let the drop, reorder and replay through. The earlier raw-frame experiment,
which hashed uncompressed YUV and signed every frame of 4K video for twenty minutes across ten configurations without a late frame or thermal throttling, remains in the supplement
(Table~\ref{tab:video}) as the computational upper bound.

\subsection{A verifier that reads its trust store from flash}
\label{sec:edge-verify}

A gateway or display device verifies against a local trust store. On the
Pi~5 we verified manifests from 500 signers with leaf and intermediate
embedded, and by thumbprint from a store held in memory and from the same
store after the page cache was dropped, so that every certificate came off
the SD card (Table~\ref{tab:edge} in the supplement). Once the certificates
are in memory, embedded and referenced verification cost the same to within
\NedgeWarmDiff\,ms, so carriage stays a question of bytes. A cold store is not
free on flash: reading the signer's certificates from the card adds
\NedgeColdEd\,ms to Ed25519 verification and \NedgeColdDSAsixtyfive\,ms to
ML-DSA-65.

\subsection{Credentials on a microcontroller}
\label{sec:mcu}

We signed and verified the same manifests on three microcontrollers and
measured the energy, and the radio cost of each carriage form over Bluetooth
LE. The Raspberry Pi Pico (RP2040, Cortex-M0+ at 125\,MHz, 264\,KiB of SRAM)
runs PQClean's portable C; the NUCLEO-L4R5ZI (STM32L4R5, Cortex-M4F,
640\,KiB) runs pqm4's assembly at 80\,MHz and at the 20\,MHz of pqm4's
benchmarks; the nRF52840-DK (Cortex-M4F at 64\,MHz, 256\,KiB) runs the same
objects and, for ES256 and SHA-256, its CryptoCell-310 through PSA Crypto.
Ed25519 and ES256 use Monocypher and micro-ecc. Each board holds one
manifest per algorithm in each carriage form; to verify, the firmware copies
it into a receive buffer as if it had just arrived, decodes it, rebuilds the
signed structure, checks the certificates up to the anchor in its trust store
and checks the manifest signature. Energy comes from a Nordic PPK2 sampling at
100\,kHz around marked batches of operations.

\begin{table}[t]
\caption{Signing and verifying on three microcontrollers, milliseconds, medians. \emph{Verify} is a complete manifest: receive it, decode it, check the leaf and intermediate up to the root anchor in the trust store, and check the manifest signature, with leaf and intermediate carried (\emph{L+I}) or only a thumbprint (\emph{Th.}). Pico: RP2040 Cortex-M0+ at 125\,MHz, portable C (PQClean). NUCLEO: STM32L4R5 Cortex-M4F at 80\,MHz, pqm4 assembly. nRF: nRF52840 Cortex-M4F at 64\,MHz, the same objects; \emph{ES256 (hw)} runs on its CryptoCell-310. Every board rejected a flipped signature bit, a tampered manifest, a tampered intermediate and a missing intermediate for every algorithm.}
\label{tab:mcu}
\small
\setlength{\tabcolsep}{4pt}
\begin{tabular}{lrrrrrrrr}
\toprule
 & \multicolumn{3}{c}{Pico} & \multicolumn{3}{c}{NUCLEO, 80\,MHz} & \multicolumn{2}{c}{nRF52840}\\
\cmidrule(lr){2-4}\cmidrule(lr){5-7}\cmidrule(lr){8-9}
Algorithm & Sign & L+I & Th. & Sign & L+I & Th. & Sign & L+I\\
\midrule
Ed25519 & 57.7 & 354.7 & 350.1 & 21.5 & 107.8 & 104.8 & 21.9 & 101.2\\
ES256 & 131.6 & 440.4 & 433.8 & 131.0 & 438.0 & 431.6 & 152.0 & 503.7\\
ES256 (hw) & -- & -- & -- & -- & -- & -- & 30.2 & 66.2\\
ML-DSA-44 & 154.8 & 160.7 & 143.3 & 96.3 & 130.9 & 112.1 & 90.5 & 121.4\\
ML-DSA-65 & 222.0 & 255.9 & 232.2 & 154.9 & 214.9 & 189.2 & 144.7 & 200.6\\
ML-DSA-87 & 329.7 & 412.5 & 381.5 & 190.0 & 360.1 & 325.7 & 199.6 & 333.1\\
\bottomrule
\end{tabular}
\end{table}


\paragraph{Time fits the budget}
On the slowest part, ML-DSA-44 and ML-DSA-65 verify faster than the classical schemes they replace, and ML-DSA-87 takes \NpicoChainDSAeightyseven\,ms (Table~\ref{tab:mcu}). The Pico checks an ML-DSA-44 manifest in
\NpicoChainDSAfortyfour\,ms and an ML-DSA-65 manifest in
\NpicoChainDSAsixtyfive, against \NpicoChainEd{} for Ed25519 and
\NpicoChainES{} for ES256. That comparison is between implementations as much
as algorithms, in the sense of Section~\ref{sec:backends}: Monocypher and
micro-ecc are compact portable libraries, and assembly-optimised Ed25519 for
the Cortex-M0 is faster than either~\cite{li2024m0}. At 20\,MHz the NUCLEO's
ML-DSA-65 signs in a median of \NnucTwentySignMcycles{} million cycles, in
line with pqm4's published counts~\cite{pqm4}. Signing takes
\NnucSignDSAsixtyfive\,ms on the NUCLEO at 80\,MHz and \NnrfSignDSAsixtyfive\,ms
on the nRF52840 for ML-DSA-65, a few frame periods at 30 frames per second,
and \NnucSignStackDSAsixtyfive\,KiB of stack.

Hardware support reverses the order. The CryptoCell-310 verifies an ES256
manifest in \NnrfChainEShw\,ms against \NnrfChainES{} in software on the same
chip, \NhwSpeedup{} times faster, which puts ES256 ahead of every ML-DSA
level: ML-DSA-44 takes \NhwRatioDSAfortyfour{} times as long and ML-DSA-87
\NhwRatioDSAeightyseven. The engine has no lattice support, so a product built
on this part gives up its accelerator when it gives up ECDSA.

\begin{table*}[t]
\caption{What a manifest costs a microcontroller beyond time. \emph{RAM}: received manifest, rebuilt signed structure and peak stack, all live at once, by carriage form (root, leaf and intermediate, leaf only, thumbprint); M0+ is the Pico's portable C, M4 the NUCLEO's pqm4 build. \emph{Verify}: energy for one complete verification, with leaf and intermediate carried, measured with the current vectors. The three boards are measured at different points (the whole Pico, the NUCLEO's MCU rail, the nRF52840's supply) and are not comparable with one another. Radio energy, measured, is in Table~\ref{tab:radio}.}
\label{tab:mcucost}
\small
\setlength{\tabcolsep}{3.6pt}
\begin{tabular}{lrrrrrrrr}
\toprule
 & \multicolumn{5}{c}{RAM (KiB)} & \multicolumn{3}{c}{Verify (mJ)}\\
\cmidrule(lr){2-6}\cmidrule(lr){7-9}
 & M0+ & \multicolumn{4}{c}{M4} & & & \\
\cmidrule(lr){3-6}
Algorithm & L+I & Root & L+I & Leaf & Th. & Pico & NUCLEO & nRF\\
\midrule
Ed25519 & 5.5 & 6.2 & 5.3 & 3.9 & 2.6 & 36.4 & 6.48 & 1.35\\
ES256 & 4.7 & 5.7 & 4.7 & 3.2 & 1.8 & 44.9 & 24.95 & 5.82\\
ES256 (hw) & -- & 9.9 & 8.9 & 7.4 & 6.0 & -- & -- & 0.89\\
ML-DSA-44 & 55.4 & 37.0 & 28.9 & 20.4 & 12.0 & 16.3 & 7.21 & 1.81\\
ML-DSA-65 & 83.3 & 47.8 & 36.8 & 25.3 & 13.8 & 25.9 & 11.84 & 3.01\\
ML-DSA-87 & 126.5 & 62.6 & 47.7 & 32.4 & 17.0 & 41.6 & 19.82 & 5.00\\
\bottomrule
\end{tabular}
\end{table*}


\paragraph{Memory and the radio}
Memory is where post-quantum credentials press on these parts
(Table~\ref{tab:mcucost} and Figure~\ref{fig:mcu}). On the NUCLEO, verifying
an ML-DSA-87 manifest needs \NramRootDSAeightyseven\,KiB with the root
carried, \NramChainDSAeightyseven\,KiB with leaf and intermediate,
\NramLeafDSAeightyseven\,KiB with the leaf alone and
\NramThumbDSAeightyseven\,KiB with a thumbprint; pqm4's verifier needs about
12\,KiB of stack, so the received and rebuilt buffers are the largest item.
The Pico's portable verifier needs \NramPicoDSAeightyseven\,KiB of its 264. On a
Cortex-M4 the carriage form, not the parameter set, decides whether a
verifier fits beside an application.

\begin{table}[t]
\caption{Measured Bluetooth LE energy for one manifest's COSE signature structure, with the certificates of each carriage form, between two nRF52840-DKs (2M PHY, 251-byte data length, 247-byte ATT MTU, 15\,ms connection interval, boards about a metre apart), on the measured board's supply, microjoules, medians of three rounds of ten transfers. \emph{Send}: from the first notification queued to the last acknowledged. \emph{Receive}: the same board as receiver, from the first notification received to the last. Both include the connected link's idle current and any link-layer retransmissions inside the window, not connection set-up or advertising. The receiving board checked every transfer byte for byte against its own copy; all arrived intact.}
\label{tab:radio}
\small
\setlength{\tabcolsep}{4pt}
\begin{tabular}{lrrrrrrrr}
\toprule
 & \multicolumn{4}{c}{Send} & \multicolumn{4}{c}{Receive}\\
\cmidrule(lr){2-5}\cmidrule(lr){6-9}
Algorithm & Root & L+I & Leaf & Th. & Root & L+I & Leaf & Th.\\
\midrule
Ed25519 & 205 & 164 & 87 & 14 & 214 & 169 & 89 & 14\\
ES256 & 228 & 173 & 91 & 13 & 242 & 180 & 94 & 14\\
ML-DSA-44 & 1,684 & 1,227 & 745 & 273 & 1,742 & 1,266 & 775 & 282\\
ML-DSA-65 & 2,310 & 1,669 & 1,017 & 366 & 2,387 & 1,713 & 1,045 & 380\\
ML-DSA-87 & 3,119 & 2,266 & 1,381 & 516 & 3,194 & 2,355 & 1,416 & 530\\
\bottomrule
\end{tabular}
\end{table}


We measured the radio both ways between two nRF52840-DKs, each manifest form
sent ten times per round in three rounds, with the measured board as sender
and then as receiver (Table~\ref{tab:radio}). Energy grows in a straight line
with the bytes: \NbleTxNj\,nJ a byte to send and \NbleRxNj\,nJ a byte to
receive (95\% intervals \NbleTxNjLo--\NbleTxNjHi{} and
\NbleRxNjLo--\NbleRxNjHi), fitted over \NbleTxN{} measured transfers from
\NbleTxXmin{} to \NbleTxXmax{} bytes with $R^2$ of \NbleTxRsq; the largest residual is \NbleTxResUj\,$\mu$J sending and \NbleRxResUj\,$\mu$J receiving, large only relative to the smallest, thumbprint-sized transfers. The windows cover the connection events
of a transfer, link-layer acknowledgements and any retransmissions, and the
connected link's idle current between them, but not connection set-up or
advertising. The radio pays for the carriage form directly. A receiver spends
\NrecvMChainDSAeightyseven\,mJ on an ML-DSA-87 manifest with leaf and
intermediate, \NrecvMLeafDSAeightyseven\,mJ on the leaf alone and
\NrecvMThumbDSAeightyseven\,mJ on a thumbprint, and receiving the chained
form costs \NrecvShareVerify\% of what verifying it costs on the same chip, so
a battery-powered verifier pays for the chain twice, once to receive it and
once to check it.

\begin{figure}[t]
\centering
\includegraphics[width=\linewidth]{figures/mcu}
\Description{Two grouped bar charts by algorithm. Left: RAM needed to verify
on the NUCLEO for four carriage forms; the ML-DSA bars fall from about 60 KiB
with the root carried to under 20 KiB with a thumbprint. Right: measured energy for an nRF52840 to receive
the manifest over Bluetooth LE in the same forms; ML-DSA-87 falls from about
3 mJ to about half a millijoule.}
\caption{Carriage form against verifier RAM on the NUCLEO-L4R5ZI
(Table~\ref{tab:mcucost}) and against the measured energy an nRF52840 spends
receiving the manifest (Table~\ref{tab:radio}).}
\label{fig:mcu}
\end{figure}

\subsection{Secure elements}
\label{sec:constrained}

The devices signing the most photographs are phones, which sign inside secure
elements, and no development board reproduces that. The measurements bound it
from both sides: an application processor signs every frame with time to
spare, and a Cortex-M4 needs a few frame periods and tens of kilobytes of
stack per ML-DSA signature. ML-DSA offers a way to split the work. The
external-$\mu$ variant lets the application processor hash the message
together with the public key into the 64-byte $\mu$ that ML-DSA signs, so a
secure element receives a fixed-size input whatever the message, and the
certificate profile allows it~\cite{fips204,rfc9881}. In a C2PA manifest the
signed message includes the protected header, and with it the embedded chain;
external $\mu$ keeps the chain's bytes out of the secure element, which then
does only the signing arithmetic instead of hashing tens of kilobytes streamed
to it over a slow bus. Whether a secure element keeps
up with burst capture remains to be measured.
